\documentclass[12pt]{article}
\usepackage[margin=0.85in]{geometry}
\usepackage{enumitem}
\usepackage{listings}
\usepackage{xcolor}
\usepackage{framed}
\usepackage{amssymb}

\lstset{
  basicstyle=\ttfamily\small,
  breaklines=true,
  frame=single,
  backgroundcolor=\color{gray!8}
}

\setlist[itemize]{leftmargin=*, itemsep=0.4em}

\title{CS 449 -- Lecture 2 Note Template}
\author{Dr.\ Jack Lange}
\date{}

\begin{document}
\maketitle

\section{Turning work in with \texttt{tar}}

This lecture starts with a lab tool, then spends most of its time on a bigger idea: how a C program gets data in and sends data out.

\begin{itemize}
  \item When would you create a \texttt{.tar.gz} versus extract one? What can go wrong if you mix up the flags?
  \vspace{2.0em}
  \item Write a command you would actually use to bundle this week's lab files, then the command a TA would use to unpack it.
  \vspace{1.6em}
\end{itemize}

\begin{lstlisting}
# create:   tar _______________________________
# extract:  tar _______________________________
\end{lstlisting}

\vspace{1.4em}

\section{Why C cares about order}

The slides contrast C with Java before talking about input. That is the setup: if you do not know how C is parsed, the later code examples will look like magic failures.

\begin{itemize}
  \item In your own words, what does ``C is parsed top down'' change about how you write a file?
  \vspace{2.2em}
  \item Compare C and Java: which kinds of names still have to appear above their use in Java, and which do not? Why might C be stricter?
  \vspace{2.2em}
  \item What would break if you called a helper function that is defined later in the same \texttt{.c} file, with no declaration above \texttt{main}?
  \vspace{1.8em}
\end{itemize}

\begin{lstlisting}
/* Predict: does this compile in C? Why or why not? */
int main(void) {
    foo();
}

void foo(void) {
    /* what would you add or move? */
}
\end{lstlisting}

\vspace{1.8em}

\section{Three ways a C program gets input}

The rest of the lecture is really one question: where can bytes come from, and who is in charge of them?

\subsection{Command-line arguments}
\begin{itemize}
  \item Imagine you run \texttt{./cmdargs hello 42}. What is in \texttt{argc}, \texttt{argv[0]}, and \texttt{argv[1]}? What is \emph{not} in \texttt{argv}?
  \vspace{2.2em}
  \item When is \texttt{argv} a better input tool than typing into the program after it starts?
  \vspace{2.0em}
\end{itemize}

\begin{lstlisting}
int main(int argc, char *argv[]) {
    // If the user forgot an argument, what should you check?
    // ________________________________
}
\end{lstlisting}

\vspace{1.2em}

\subsection{Console input with \texttt{scanf}}
\begin{itemize}
  \item \texttt{scanf} does not just ``read text.'' What job is the format string doing, and why does the \texttt{\&} matter?
  \vspace{2.2em}
  \item A classmate writes \texttt{scanf("\%d", n);} (no \texttt{\&}) and the program crashes or silently corrupts memory. What did they misunderstand?
  \vspace{2.2em}
  \item Why does the lecture care about the return value? Invent a tiny input that should make \texttt{scanf} fail.
  \vspace{2.0em}
\end{itemize}

\begin{lstlisting}
int n;
int matched = scanf("%d", /* what belongs here? */);
/* If matched is 0, what happened, and what is in n? */
\end{lstlisting}

\vspace{1.8em}

\section{Streams: the idea that ties input and output together}

\texttt{argv} is a one-shot list. \texttt{scanf} is already talking to a stream. The lecture's third mechanism is the general version of that idea.

\begin{itemize}
  \item Explain FIFO like you are teaching a friend who has never heard the term. Why does ``in order'' matter for a console?
  \vspace{2.2em}
  \item What is each of \texttt{stdin}, \texttt{stdout}, and \texttt{stderr} \emph{for}? When would you print an error to \texttt{stderr} instead of \texttt{stdout}?
  \vspace{2.4em}
  \item The slides also mention files, devices, and network sockets as streams. What is the shared mental model?
  \vspace{2.0em}
\end{itemize}

\noindent\textbf{Diagram: sketch a stream from your program, through the OS, to the console (or a file). Mark the buffer and a flush.}
\begin{framed}
\begin{minipage}[t][0.35\textheight]{\textwidth}
\vspace{0.6em}
\textit{Do not copy a figure from memory. Draw the pieces and label who waits if the buffer is not flushed.}
\end{minipage}
\end{framed}

\vspace{1.2em}

\subsection{Buffering, \texttt{printf}, and \texttt{fprintf}}
\begin{itemize}
  \item Why can a program look ``stuck'' or print nothing until it exits, even if you called \texttt{printf}?
  \vspace{2.2em}
  \item \texttt{printf} picks a stream for you. \texttt{fprintf} makes you pick. When is the extra argument worth it?
  \vspace{2.2em}
\end{itemize}

\begin{lstlisting}
int printf(const char *format, ...);
int fprintf(FILE *stream, const char *format, ...);
/* Write one situation where fprintf is the better call. */
/* ___________________________________________________ */
\end{lstlisting}

\vspace{1.8em}

\section{The shell already knows about streams}

Unix ``everything is a file'' is why the same program can talk to a keyboard, a file, or another program without being rewritten.

\begin{itemize}
  \item In your own words, why can the shell mix files and console streams so easily?
  \vspace{2.2em}
  \item You ran \texttt{ls > files.out} yesterday and today you run \texttt{ls >> files.out}. What is different about the file afterward? When is each the one you want?
  \vspace{2.4em}
  \item Predict what \texttt{wc < files.out} does without looking it up from the slide. Who is ``typing'' into \texttt{wc}?
  \vspace{2.2em}
  \item Walk through \texttt{ls | grep .c | wc} as a story: what does each stage receive, and what does it pass on?
  \vspace{2.4em}
\end{itemize}

\vspace{0.6em}

\section{Special streams the OS leaves lying around}

These look like files in the filesystem, but they are features.

\begin{itemize}
  \item A program is too noisy. You want to keep the program and throw away its output. Which special file do you aim at, and what happens to reads versus writes?
  \vspace{2.4em}
  \item The slides say \texttt{/dev/zero} is different from \texttt{/dev/null}. How?
  \vspace{2.4em}
  \item When might a program read from \texttt{/dev/random} instead of those two?
  \vspace{2.0em}
\end{itemize}

\vspace{1.2em}

\section{When you, not the shell, open the file}

Command-line arguments and the three console streams are set up for you. A file you name in code is your problem.

\begin{itemize}
  \item Why can a text file, once opened, be treated like the console? What functions become the file versions of \texttt{printf} / \texttt{scanf}?
  \vspace{2.2em}
  \item What does a \texttt{NULL} from \texttt{fopen} mean in real life (permissions, bad path, \ldots)? What should a careful program do next?
  \vspace{2.4em}
\end{itemize}

\begin{lstlisting}
FILE *file = fopen(name, /* which mode, and why? */);
if (file == NULL) {
    /* what belongs here? */
}
/* talk to the file */
fclose(file);
\end{lstlisting}

\vspace{1.8em}

\section{What actually lands on disk: ASCII versus raw bytes}

This is the lecture's ``aha'' about why printing a number and writing a number are not the same.

\begin{itemize}
  \item What problem is ASCII solving? Why is the character \texttt{"7"} not the same thing as the integer 7 sitting in a variable?
  \vspace{2.4em}
  \item The slides show \texttt{fprintf("\%d\textbackslash n", i)} for $i = 7$ and $i = -7$. Before you copy their hex, predict: why do extra bytes appear in the file that were not the in-memory integer?
  \vspace{2.4em}
\end{itemize}

\noindent\textbf{Diagram: for one of those examples, show memory on the left and the file on the right. Label which side is ``the number'' and which side is ``text.''}
\begin{framed}
\begin{minipage}[t][0.35\textheight]{\textwidth}
\vspace{0.6em}
\textit{Use the slide values if you want ($0$x$00000007$ vs $0$x$37$, $0$x$0$a). The point is the comparison, not the hex trivia.}
\end{minipage}
\end{framed}

\vspace{1.2em}

\section{Explicit copies: \texttt{fread} and \texttt{fwrite}}

\begin{itemize}
  \item The lecture says an application can only operate on data in memory. Why does that force you to copy to and from files yourself?
  \vspace{2.2em}
  \item File access is ``like memory, but explicit.'' What is implicit about reading a variable, and what is explicit about \texttt{fread}?
  \vspace{2.4em}
  \item In \texttt{fwrite(\&x, sizeof(x), 1, outfile)}, what is each argument responsible for? What goes wrong if you pass \texttt{x} instead of \texttt{\&x}?
  \vspace{2.4em}
\end{itemize}

\begin{lstlisting}
int x = 5;
FILE *outfile = fopen(argv[1], "wb");
fwrite(&x, sizeof(x), 1, outfile);
/* After this, how does the file differ from fprintf("%d", x)? */
fclose(outfile);
\end{lstlisting}

\vspace{1.8em}

\section{The file position you cannot see}

\texttt{fread} / \texttt{fwrite} do not take an offset. Something else is remembering where you are.

\begin{itemize}
  \item If a file is an array of bytes, who is keeping the current index, and when does that index move?
  \vspace{2.2em}
  \item You read 5 bytes, then read again. Where does the second read start? Why can that surprise you?
  \vspace{2.2em}
  \item What is \texttt{fseek(fstream, 16, SEEK\_SET)} asking the system to do? Invent a reason you would need that instead of reading from the start every time.
  \vspace{2.4em}
\end{itemize}

\noindent\textbf{Diagram: draw a file as a row of bytes. Mark the position before a 5-byte read, after it, and after that \texttt{fseek}.}
\begin{framed}
\begin{minipage}[t][0.35\textheight]{\textwidth}
\vspace{0.6em}
\textit{Leave room for arrows. The picture should explain the API, not decorate the page.}
\end{minipage}
\end{framed}

\vspace{2.2em}
\begin{center}
\textit{This template may not cover every topic from the source slides.}
\end{center}

\end{document}
